% titel: Blatt 1, A1: Negation einer Aussage

% Teil I: Werkzeuge, kein Lösungsweg.

Von einer Aussage das Gegenteil (die Negation) unter den vorgegebenen
Möglichkeiten erkennen; nicht über den Wahrheitsgehalt entscheiden.

\begin{werkzeugkasten}[Negation]
Die Negation einer Aussage mit Allquantor ``zu jeder Zahl gibt es eine
natürliche Zahl, die größer ist'' ist die Aussage: \emph{Es gibt eine Zahl,
die von keiner natürlichen Zahl übertroffen wird.}
Formal:
\begin{equation}
  \neg\bigl(\forall x\,\exists n\in\mathbb N:\ n>x\bigr)\;\Longleftrightarrow\;
  \exists x\,\forall n\in\mathbb N:\ n\leq x .
\end{equation}
``Zu jeder'' wird zu ``es gibt eine'', ``es gibt eine'' wird zu ``zu keiner''.
\end{werkzeugkasten}

\begin{achtung}[Darauf achten]
\begin{itemize}
  \item Die Negation von ``es gibt eine natürliche Zahl, die größer ist''
        ist ``es gibt keine natürliche Zahl, die größer ist'', nicht
        ``kleiner''.
  \item Aussagen wie ``Keine einzige Zahl wird von jeder natürlichen Zahl
        übertroffen'' (Variante a) sind \emph{nicht} die Negation, weil sie
        die Reihenfolge der Quantoren falsch drehen.
  \item Nicht bewerten ob die Aussage wahr oder falsch ist.
\end{itemize}
\end{achtung}
